

\documentclass[11pt,a4paper]{article}
\usepackage[utf8]{inputenc}
\usepackage[T1]{fontenc}
\usepackage[slovak]{babel}
\usepackage{geometry}
\usepackage{tabularx}
\usepackage{longtable}
\usepackage{array}
\usepackage{xcolor}
\usepackage{booktabs}
\usepackage{hyperref}
\usepackage{enumitem}
\geometry{margin=2.2cm}
\setlist[itemize]{leftmargin=1.2em, topsep=3pt, itemsep=1pt}
\setlist[enumerate]{leftmargin=1.5em, topsep=3pt, itemsep=1pt}

\definecolor{apvvgray}{RGB}{240,240,240}
\newcolumntype{Y}{>{\raggedright\arraybackslash}X}
\newcommand{\Field}[2]{\rowcolor{apvvgray}\textbf{#1} & #2\\}


\begin{document}
\begin{center}
{\Large \textbf{VV–F – Vecný zámer projektu}}\\[2pt]
\textit{Záväzná osnova pre aplikovaný výskum a vývoj}
\end{center}

\noindent\textbf{Názov projektu:} TimeManager: Inteligentný plánovač času pre študentov s integráciou Google služieb (TIMA)\\
\textbf{Zodpovedný riešiteľ:} Michal Pípa \quad
\textbf{Žiadateľ:} Slovenská technická univerzita v Bratislave, FIIT STU \quad
\textbf{Štatutárny zástupca:} \emph{doplní žiadateľ}

\section{Excelentnosť}
\subsection{Projektový zámer, spôsob riešenia a aplikačná úroveň}
TIMA je webová aplikácia pre študentov s AI asistenciou. Využíva Next.js 14 (App Router, RSC) a Django REST API s PostgreSQL (JSONB, full-text).
Kľúčom je AI prioritizácia (Gemini) a obojsmerná synchronizácia s Google Calendar/Tasks.
Aplikačná úroveň: softvérová služba s pilotnou prevádzkou na univerzite, pripravená na širší rollout.

\subsection{Aktuálnosť a originalita}
\begin{itemize}
\item Reaguje na fragmentáciu nástrojov (kalendár, úlohy, poznámky) v študentskom prostredí.
\item Kombinuje Eisenhowerovu maticu, Pomodoro a AI do jedného nástroja.
\item Kontextové odporúčania na základe záťaže a deadline.
\end{itemize}

\subsection{Stav poznania a predchádzajúce výsledky}
\begin{itemize}
\item Analýza metód time managementu (Eisenhower, Pomodoro).
\item Benchmarky existujúcich riešení (Todoist, Notion, Google Calendar).
\end{itemize}

\subsection{Ciele a reálnosť ich dosiahnutia}
\begin{enumerate}
\item Navrhnúť a implementovať škálovateľné jadro (API, databáza) do 12 mesiacov.
\item Integrovať AI modul pre prioritizáciu a odhady času do 18 mesiacov.
\item Dodať UX s dashboardom, Pomodoro a tímovou spoluprácou do 24 mesiacov.
\item Pilotná prevádzka na FIIT STU a vyhodnotenie dopadu do 36 mesiacov.
\end{enumerate}

\subsection{Metodológia riešenia}
\begin{itemize}
\item Iteratívny vývoj (agile), CI/CD, jednotkové a integračné testy.
\item Zber metrik produktivity a SUS dotazník pri pilote.
\item Bezpečnostný dizajn (OAuth2, JWT, šifrovanie, GDPR).
\end{itemize}

\subsection{Najdôležitejšie výstupy riešiteľa (posledných 5 rokov)} 
\begin{itemize}
\item Prototyp TIMA – otvorené moduly (UI komponenty, synchronizačný modul).
\item Odborný článok o architektúre a integráciách (študentské fórum).
\item Interné technické správy (databázová schéma, API špecifikácie).
\end{itemize}

\subsection{Najdôležitejšie projekty riešiteľa (posledných 5 rokov)}
Študentské projekty zamerané na plánovanie, integrácie API, UI/UX (FIIT STU).

\subsection{Kvalita a vízia}
\begin{itemize}
\item Budovanie tímu (backend, frontend, integrácie, UX).
\item Udržateľnosť – freemium model pre školy, otvorené moduly.
\end{itemize}

\subsection{Kompetentnosť organizácií a tímu a zapojenie mladých výskumníkov}
FIIT STU disponuje zázemím pre softvérový výskum; tím pokrýva backend, frontend, integrácie a UX.
Projekt je vedený mladým riešiteľom a zapája študentov magisterského a PhD. štúdia.

\section{Dopad}
\begin{itemize}
\item Nová služba pre študentov so zlepšenou produktivitou a wellbeingom.
\item Overené postupy – AI prioritizácia, obojsmerná synchronizácia, vizuálne plánovanie.
\item Využiteľnosť na univerzitách v SR a v zahraničí (multijazyčnosť).
\item Ekonomické prínosy: zníženie času na koordináciu, lepšia úspešnosť štúdia.
\item Popularizácia: web, workshopy, akademické podujatia, open-source moduly.
\end{itemize}

\noindent\textbf{Opatrenia na maximalizáciu dopadu:}
\begin{enumerate}
\item Pilot s min. 200 používateľmi, priebežné meranie metrik.
\item Štandardizovaná integrácia s LMS (Moodle) a Microsoft Teams.
\item Prístupnosť a inklúzia (WCAG) – širšia použiteľnosť.
\end{enumerate}

\section{Implementácia}
\subsection{Časový harmonogram (01.01.2026 – 31.12.2028) a pracovné balíky}
\begin{itemize}
\item \textbf{WP1: Analýza a návrh} (01/2026 – 06/2026): Špecifikácia požiadaviek; architektúra; GDPR DPIA draft.
\item \textbf{WP2: Backend \& API} (04/2026 – 12/2026): API v1; AI modul v1; Sync v1.
\item \textbf{WP3: Frontend \& UX} (07/2026 – 06/2027): UI prototyp; Beta UI; UX testy v1.
\item \textbf{WP4: Pilot a evaluácia} (07/2027 – 12/2028): Pilot v1; Report metrik v1; Záverečná evaluácia.
\item \textbf{WP5: Disseminácia a udržateľnosť} (01/2026 – 12/2028): Web \& dokumentácia; workshopy; plán udržateľnosti.
\end{itemize}

\subsection{Projektový manažment}
SCRUM (2–3 týždňové sprinty), mesačné revízie, rizikový register.

\subsection{Riziká a mitigácia}
\begin{itemize}
\item Rate limiting Google API $\rightarrow$ batchovanie, caching, exponenciálne backoff.
\item GDPR a bezpečnosť $\rightarrow$ DPIA, šifrovanie, minimizácia dát.
\item Náklady na AI $\rightarrow$ prompt optimalizácia, cache, plán B s lokálnymi modelmi.
\end{itemize}

\subsection{Rozpočet a zdroje}
Rozpočet je uvedený v časti VV–C: mzdové a osobné náklady; služby (cloud, AI, bezpečnostný audit, UX výskum); materiál; cestovné; popularizácia; energie/komunikácie; nepriame náklady.

\subsection{Technická a personálna infraštruktúra}
Vývojové servery, repozitáre, CI/CD; skúsenosti s Next.js, Django, PostgreSQL, Google API.

\vfill
\noindent\textit{Pozn.: Obsah je zosúladený s \textbf{VV–F} štruktúrou a údajmi uvedenými v \textbf{VV–A až VV–E}.}
\end{document}
